% !TeX root = ../main.tex
\section{Тестовое воздействие и результаты имитационного моделирования}

Для проверки работоспособности разработанного тракта цифровой обработки используется тестовое воздействие в виде симметричного линейно-частотно модулированного (ЛЧМ) радиоимпульса с прямоугольной огибающей.Сигнал генерируется с базой $B = 512$ и шириной спектра $\Delta f = 20$ МГц, имитируя прием на фоне белого гауссова шума. На рисунке 2 представлен фрагмент зашумленного высокочастотного сигнала на выходе аналого-цифрового преобразователя перед началом цифровой обработки.

\begin{figure}[H]
    \centering
    \includegraphics[width=0.9\textwidth]{1.png}
    \caption{Сигнал на выходе АЦП (фрагмент)}
    \label{fig:adc_out}
\end{figure}

После выполнения этапов цифровой демодуляции (перемножения с цифровым гетеродином), фильтрации нижних частот и децимации, спектр сигнала переносится в область нулевой частоты. На рисунке 3 показана сформированная комплексная огибающая (синфазная $I$ и квадратурная $Q$ составляющие) сигнала на выходе тракта предварительной обработки.

\begin{figure}[h!]
    \centering
    \includegraphics[width=0.9\textwidth]{2.png}
    \caption{Комплексная огибающая после ЦДМ, ЦФ и децимации}
    \label{fig:baseband_out}
\end{figure}

Окончательная внутриимпульсная обработка производится в согласованном фильтре с использованием алгоритма сегментированной быстрой линейной свертки (метод перекрытия со сложением). На рисунке 4 представлен отклик фильтра --- сжатый по времени импульс. Сравнение двух методов показывает, что применение весового окна Хэмминга к опорному сигналу позволяет существенно снизить уровень боковых лепестков сжатого сигнала по сравнению с прямоугольным окном (без взвешивания), хотя и приводит к закономерному, но незначительному расширению главного лепестка.

\begin{figure}[H]
    \centering
    \includegraphics[width=0.9\textwidth]{3.png}
    \caption{Выход согласованного фильтра}
    \label{fig:mf_out}
\end{figure}